\chapter{Conclusões e trabalhos futuros}\label{sec:conclusoes}

A realização do estágio supervisionado possibilitou o cumprimento dos objetivos propostos,
com o desenvolvimento de um sistema web de agendamento de salas para a UniFil.
A aplicação substituiu o controle feito em planilhas e estabeleceu um fluxo único
de solicitação, análise e confirmação de reservas.

Ao longo do processo foi possível aplicar na prática conceitos estudados em sala de aula,
como o levantamento de requisitos, o mapeamento de processos em BPMN e a modelagem de sistemas.
O maior desafio esteve em abstrair e gerenciar a complexidade do software,
definindo as fronteiras entre os módulos e separando a regra de negócio
do detalhe de tecnologia, o que levou a organizar a aplicação em camadas
com as dependências voltadas para o domínio.

Como possibilidade de estudos futuros, destaca-se o envio automático de notificações
sobre mudanças no status das solicitações e a geração de um relatório de grade de horários,
consolidando a ocupação das salas ao longo do semestre.
Sugere-se ainda uma funcionalidade de solicitação de troca de salas entre professores,
trazendo para dentro do sistema a negociação que hoje ocorre de maneira informal.
Mostra-se relevante também a integração com um servidor MCP (\textit{Model Context Protocol}),
permitindo que assistentes de inteligência artificial consultem a disponibilidade
e solicitem reservas.
